\subsubsection{高阶方程与方程组}\label{subsec:higher-order-systems}

将上述定义推广到二阶以上的高阶方程是直接的.

\begin{Definition}
\label{def:higher-order-pde-type}
设 $L$ 为 \eqref{eq:general-linear-differential-expression} 中定义的 $m$ 阶算子. 特征曲面由方程
\begin{equation}\label{eq:higher-order-characteristic-surface}
L^p(\bm{x},\nabla\phi)=0.
\end{equation}
定义. 如果方程在 $\bm{x}$ 处没有实特征, 或等价地, 在 $\bm{x}$ 处有
\begin{equation}\label{eq:higher-order-ellipticity}
L^p(\bm{x},i\bm{\xi})\ne0,\qquad\forall\bm{\xi}\ne0.
\end{equation}
则称它在该处是\textbf{椭圆型的}. 若满足以下条件, 则称方程沿方向 $\bm{n}$ 是\textbf{严格双曲型的}\footnote{虽然存在非严格双曲型的一般定义, 但我们不在这里给出. 下文将为一阶方程组定义非严格双曲性.}:
\begin{enumerate}
\item $L^p(\bm{x},i\bm{n})\ne0$, 并且
\item 对每个不与 $\bm{n}$ 共线的 $\bm{\xi}\in\mathbb{R}^n$, 方程
\begin{equation}\label{eq:strict-hyperbolic-characteristic-roots}
L^p(\bm{x},i\bm{\xi}+i\omega\bm{n})=0
\end{equation}
的所有根 $\omega$ 都为互异的实数.\end{enumerate}
\end{Definition}

在应用中, $\bm{n}$ 通常是与时间相联系的坐标方向. 此时, 令 $\bm{x}=(x_1,x_2,\ldots,x_{n-1},t)$, 并令 $\bm{\xi}=(\xi_1,\ldots,\xi_{n-1},0)$ 为空间向量.

对于支集很小且快速振荡的函数, 可以将 $L^p$ 的系数看作近似常数; 现在假设它们就是常数. 如果 $\omega$ 是 \eqref{eq:strict-hyperbolic-characteristic-roots} 的根, 那么 $u=\exp(i(\bm{\xi}\cdot\bm{x})+i\omega t)$ 是 $L^pu=0$ 的一个解. 如果 $\omega$ 的虚部为负, 这一解就随时间指数增长. 此外, 由于 $L^p$ 是 $m$ 次齐次的, 即对任意标量 $\lambda$ 有 $L^p(\bm{x},\lambda(i\bm{\xi}+i\omega\bm{n}))=\lambda^mL^p(\bm{x},i\bm{\xi}+i\omega\bm{n})$, 所以只要有非实根, 就总有虚部为负的根 (如果改变 $\bm{\xi}$ 的符号, 也就改变了 $\omega$ 的符号). 再者, 如果将 $\bm{\xi}$ 乘以标量因子 $\lambda$, 则 $\omega$ 也乘以同一因子, 因而解在空间振荡得越快, 其增长就会越快. 所以 \eqref{eq:strict-hyperbolic-characteristic-roots} 中的根为实数, 是初值问题适定的一个必要条件.

现在将注意力转向包含 $k$ 个未知函数 $u_j$, $j=1,2,\ldots,k$ 的 $k$ 个偏微分方程组成的方程组:
\begin{equation}\label{eq:linear-pde-system}
L_{ij}(\bm{x},D)u_j=0,\qquad i=1,2,\ldots,k.
\end{equation}
同代数方程组一样, 适定问题要求方程个数与未知函数个数相等, 所以假设算子 $L_{ij}$ 构成方阵 $\bm{L}$. 原则上, 上述概念的推广是很直接的.

\begin{Definition}
\label{def:system-pde-type}
特征曲面由方程
\begin{equation}\label{eq:system-characteristic-surface}
\det\bm{L}^p(\bm{x},\nabla\phi)=0,
\end{equation}
定义, 没有实特征曲面的方程称为\textbf{椭圆型的}. \textbf{严格双曲性}也按上述方式定义, 只需将定义中的 $L^p$ 替换为 $\det\bm{L}^p$. 若 $\bm{L}^p$ 的所有分量都是一阶算子, 则当以下条件成立时, 称方程组沿方向 $\bm{n}$ 是双曲型的 (不一定严格):
\begin{enumerate}
\item $\det\bm{L}^p(\bm{x},\bm{n})\ne0$, 并且
\item 对于不与 $\bm{n}$ 共线的 $\bm{\xi}$, 问题
\[
\det\bm{L}^p(\bm{x},i\bm{\xi}+i\omega\bm{n})=0
\]
的所有特征值 $\omega$ 都是实数, 且存在一组完备的特征向量.\end{enumerate}
\end{Definition}

注意, 由于已假设 $\bm{L}^p$ 的分量都是一阶的, 有 $\bm{L}^p(\bm{x},i\bm{\xi}+i\omega\bm{n})=\bm{L}^p(\bm{x},i\bm{\xi})+\omega\bm{L}^p(\bm{x},i\bm{n})$; 因而问题 $\det\bm{L}^p=0$ 是关于 $\omega$ 的矩阵特征值问题. 如果特征值互异, 就总存在一组完备的特征向量; 所以严格双曲性蕴含双曲性.

一般地, 定义方程组的"主部"时需要谨慎. 简单地取"最高阶项"这一朴素方法并不令人满意.

\begin{Example}
\label{ex:weighted-laplace-reformulation}
为看清这个问题, 考虑二维 Laplace 方程
\begin{equation}\label{eq:first-order-laplace-reformulation-origin}
u_{xx}+u_{yy}=0,
\end{equation}
令 $v=u_x$, $w=u_y$, 将它改写为一阶方程组. 得到的方程组为
\begin{equation}\label{eq:first-order-laplace-system}
u_x=v,\qquad u_y=w,\qquad v_x+w_y=0.
\end{equation}
如果将 $\bm{L}^p$ 定义为包含一阶项的部分, 容易看出 $\det\bm{L}^p$ 恒为零. 另一方面, 由于 Laplace 方程是椭圆型方程的标准例子, 我们希望等价的一阶方程组也被定义为"椭圆型的". 显然, 此时不能丢掉 \eqref{eq:first-order-laplace-system} 前两个方程中的项 $v$ 和 $w$.
\end{Example}

可以通过给每个方程指定"权" $s_i$, 给每个因变量指定权 $t_j$, 使各算子 $L_{ij}$ 的阶不超过 $s_i+t_j$, 来解决这一困难. 然后将主部 $L_{ij}^p$ 定义为那些阶恰好等于 $s_i+t_j$ 的项所构成的部分. 假定可以选取权, 使 $\det\bm{L}^p$ 不恒为零\footnote{有些例子不满足这一假设, 例如方程组 $u_x+v_y=0$, $u_x+v_y+v=0$. 如果使用等价形式 $u_x+v_y=0$, $v=0$, 这一困难便消失了. 已知, 对非退化方程组总存在具有所需性质的权 \cite{Vo}. 这里非退化的含义如下: 如果按通常方式将 $\det\bm{L}(i\bm{\xi})$ 表示为乘积之和, 那么其中每个乘积作为 $\bm{\xi}$ 的多项式的次数都不超过行列式的次数.}; 此时, $\det\bm{L}^p$ 由 $\det\bm{L}$ 中出现的所有 $\sum_{i,j}s_i+t_j$ 阶项构成. 例如, 在 \eqref{eq:first-order-laplace-system} 中, 可取 $s_1=s_2=t_2=t_3=0$ 和 $t_1=s_3=1$. (这里变量的次序约定为 $u,v,w$.) 采用这些权时, \eqref{eq:first-order-laplace-system} 的主部实际上与 \eqref{eq:first-order-laplace-system} 完全相同, 并且计算得到
\begin{equation}\label{eq:weighted-laplace-principal-determinant}
\det\bm{L}^p(i\bm{\xi})=\det\begin{pmatrix}i\xi_1&-1&0\\i\xi_2&0&-1\\0&i\xi_1&i\xi_2\end{pmatrix}=-\xi_1^2-\xi_2^2,
\end{equation}
它等于 Laplace 方程的符号.

对方程组使用"阶"这个术语时, 有两种不同含义. "一阶"或"二阶"方程组等术语, 通常指组成方程组的各个方程. 不过, 也可以为整个方程组指定一个阶. 这一阶定义为 $\det\bm{L}^p$ 的符号的次数, 它等于 $\sum_{i,j}s_i+t_j$.

\begin{Remark}
\label{rem:nonunique-principal-weights}
注意, 权不一定唯一. 以下是一个简单例子, 其中方程组在不止一种权的选择下都是椭圆型的. 两种选择都有用, 并会得到不同的椭圆正则性结果. 考虑方程组
\begin{equation}\label{eq:two-weight-elliptic-system}
\Delta u-v=0,\quad\Delta v=0.
\end{equation}
可以选择权 $s_1=s_2=2$, $t_1=t_2=0$. 在这些选择下, 符号的主部为
\begin{equation}\label{eq:first-choice-weighted-principal-symbol}
\bm{L}^p(i\bm{\xi})=\begin{pmatrix}-|\bm{\xi}|^2&0\\0&-|\bm{\xi}|^2\end{pmatrix}.
\end{equation}
但也可以选择权 $t_1=2$, $t_2=0$, $s_1=0$, $s_2=2$. 此时主部为
\begin{equation}\label{eq:second-choice-weighted-principal-symbol}
\bm{L}^p(i\bm{\xi})=\begin{pmatrix}-|\bm{\xi}|^2&-1\\0&-|\bm{\xi}|^2\end{pmatrix}.
\end{equation}
在这两种情形下, 都有 $\det\bm{L}^p(i\bm{\xi})=|\bm{\xi}|^4$.
\end{Remark}

\subsubsection{非线性方程}\label{subsec:nonlinear-classification}

对于非线性方程和方程组, 类型不仅可能依赖于空间中的点, 还可能依赖于解本身. 只需在给定解处将方程线性化, 并将其类型定义为线性化方程的类型. 类似地, 将特征曲面定义为线性化方程的特征曲面.

为以后使用, 给出拟线性和半线性的定义:

\begin{Definition}
\label{def:quasilinear-semilinear}
如果最高阶导数只线性地出现 (其系数可以依赖于低阶导数), 则称方程组为\textbf{拟线性的}. 如果它是拟线性的, 且最高阶项的系数只依赖于 $\bm{x}$ 而不依赖于解, 则称其为\textbf{半线性的}.
\end{Definition}

\begin{Example}
\label{ex:nonlinear-classification}
方程
\begin{equation}\label{eq:quasilinear-type-example}
\alpha(u_x)u_{xx}+u_{yy}=0
\end{equation}
是拟线性的; 若 $\alpha(u_x)>0$, 则为椭圆型, 若 $\alpha(u_x)<0$, 则为双曲型. 方程
\begin{equation}\label{eq:nonquasilinear-type-example}
\phi(u_{xx})+u_{yy}=0
\end{equation}
不是拟线性的; 若 $\phi'(u_{xx})>0$, 则为椭圆型, 若 $\phi'(u_{xx})<0$, 则为双曲型. 方程
\begin{equation}\label{eq:semilinear-elliptic-example}
(x^2+1)u_{xx}+u_{yy}+\phi(u_x,u_y)=0
\end{equation}
是半线性椭圆型方程.
\end{Example}

\paragraph*{Problems}

\begin{enumerate}[label=\thesection.\arabic*.,ref=\thesection.\arabic*,start=1,series=book-problems]

\item\label{ex:determine-pde-types} 确定下列方程的类型:\begin{enumerate}[label=(\alph*)]
\item $u_{xy}+2u_x=0$.
\item $u_{xxxx}-u_{xxyy}+u_{yyyy}=0$.
\item $u_{tt}+u_{xxxx}=0$.
\end{enumerate}

\item\label{ex:find-second-order-characteristics} 求下列方程的特征.\begin{enumerate}[label=(\alph*)]
\item $xu_{xx}+u_{yy}=0$.
\item $yu_{xx}+u_{yy}=0$.
\end{enumerate}

\item\label{ex:first-order-characteristic-solution} 考虑一阶方程 $yu_x+(x^2+1)u_y=0$.
\begin{enumerate}[label=(\alph*)]
\item 确定其特征.
\item 证明, 最一般的解是沿特征为常数的任意函数, 并利用这一事实给出通解公式.
\end{enumerate}

\item\label{ex:stokes-ellipticity} $\mathbb{R}^3$ 中的 Stokes 方程组由以下方程组成:
\begin{equation}\label{eq:classification-stokes-system}
\Delta\bm{u}-\nabla p=0,\qquad\operatorname{div}\bm{u}=0.
\end{equation}
这里的未知函数为向量函数 $\bm{u}:\mathbb{R}^3\to\mathbb{R}^3$ 和标量函数 $p:\mathbb{R}^3\to\mathbb{R}$. 证明该方程组为椭圆型.

\item\label{ex:euler-system-type} $\mathbb{R}^3$ 中的 Euler 方程为
\begin{equation}\label{eq:classification-euler-system}
(\bm{u}\cdot\nabla)\bm{u}-\nabla p=0,\qquad\operatorname{div}\bm{u}=0.
\end{equation}
这里的未知函数为向量函数 $\bm{u}:\mathbb{R}^3\to\mathbb{R}^3$ 和标量函数 $p:\mathbb{R}^3\to\mathbb{R}$. 证明该方程组既不是椭圆型, 也不是双曲型.

\item\label{ex:hyperbolicity-two-matrices} 考虑方程组
\begin{equation}\label{eq:hyperbolic-two-spatial-directions-system}
\bm{u}_t+\bm{A}\bm{u}_x+\bm{B}\bm{u}_y=0.
\end{equation}
$\bm{A}$ 和 $\bm{B}$ 必须满足什么条件, 才能使方程组沿 $t$ 方向为双曲型? 你所找到的条件一般难以验证. 能否给出一个简单特例? 现在考虑 $\bm{A}$ 和 $\bm{B}$ 都为对角矩阵的特殊情形. 在什么条件下方程组是严格双曲型的?

\end{enumerate}
